\section{Background and Related Work}

\subsection{Persistent Messaging and Stream Processing}

We consider persistent subscription transports in which a client keeps a
long-lived connection to a subscriber service and multiplexes multiple
subscriptions over that connection. The abstraction covers STOMP-over-TCP,
custom TCP subscription streams, and WebSocket-based messaging. In each case,
the load created by one connection is not captured by the connection count
alone; it also depends on how many subscriptions are multiplexed over the
connection and how much backlog has accumulated behind those subscriptions.

Modern event-driven platforms and stream-processing engines face related
elasticity problems. Dynamic stream redirection can improve runtime elasticity
in distributed stream-processing systems \cite{Liu2020DrStream}. Cardellini
et al. survey runtime-adaptation strategies and organize them by their control
objectives and mechanisms \cite{Cardellini2022DSPAdaptation}. Recent adaptive
join-processing work also shows that workload redistribution can improve
latency and throughput under changing stream rates \cite{Sun2024AjStream}.
These studies motivate adaptive runtime control. Their actions change
operator placement, stream routing, or stateful processing inside the stream
runtime. They do not address the client-side persistent-transport problem in
which subscription count and per-subscription backlog must be balanced across a
client's own long-lived connections.

Fragkoulis et al. distinguish elasticity, reconfiguration, state management,
and fault tolerance in the evolution of stream-processing systems
\cite{Fragkoulis2024StreamEvolution}. This distinction matters here because
the durable queue stays in place: migration changes its subscription's delivery
owner, not the queue's stored contents. Transactional stream processing also
separates processing guarantees from ordinary stream execution
\cite{Zhang2024TransactionalStreams}. Accordingly, a completed connection
replacement is not, by itself, evidence of end-to-end exactly-once processing.

\subsection{Autonomic and Learning-Based Control}

Gheibi et al. organize machine-learning support for self-adaptation around
monitor--analyze--plan--execute feedback loops, including policy updates and
resource management \cite{Gheibi2021MLSelfAdaptation}. Our controller follows
that established pattern; DQL alone is not the differentiating mechanism.
The industrial survey by Weyns et al. also examines the practical motivations,
risks, and difficulties of deploying self-adaptation
\cite{Weyns2023Industry}. This motivates making the observation boundary and
the concrete connection-handoff procedure explicit.

Recent DRL literature is also clear about the risks of applying learning
methods directly to production systems. Partial observability, non-stationary
workloads, reward design, and safe online adaptation remain difficult
\cite{DulacArnold2021}. Cloud DRL surveys make the same point for operational
cloud control \cite{Zhou2024DRLCloudReview}. The closer load-balancing
comparisons already include decentralized operation. The QEdgeProxy preprint
describes multi-player bandits and per-client QoS feedback for selecting
service instances \cite{Cilic2025}. Yahya et al. combine mean-field bandits
with server-reward adjustment for decentralized task offloading
\cite{Yahya2024}. BLOC uses overload feedback and dynamic capacity estimation
without a centralized metadata store \cite{Bhattacharya2024}. Selvakumar
et al. use predicted waiting times and service-chain communication costs to
prioritize requests and select instances \cite{Selvakumar2023}. These
mechanisms provide direct comparisons for decision inputs and adaptive
selection. Our narrower distinction is conditional reuse of a retained probe
socket together with a concrete subscription-ownership handoff.

Negi et al. classify soft-computing approaches to VM- and physical-machine-level
load balancing \cite{Singh2023SoftCloudLB}. In our messaging setting, a bounded
action-template design limits the corresponding control space: the DQL agent
chooses among deployable rebalancing actions rather than directly searching
the full combinatorial subscription-assignment space.

Metzger et al. use feature models to guide exploration through large or
evolving adaptation spaces \cite{Metzger2024GuidedExploration}. Our bounded
templates reduce a different search space: the possible subscription moves
and connection replacements available to one client. Recent RS-DRL work
instead changes reward shaping during replay to manage uncertainty in
self-adaptation \cite{Kavianifar2026RSDRL}. We retain a conventional DQL
learner and specialize the observations, feasible actions, and handoff
semantics to persistent messaging.

\subsection{Research Gap and Novelty}

The comparison above identifies a specific gap in how actions are chosen and
carried out. Decentralized load balancing already exists. BLOC, QEdgeProxy, and
decentralized offloading already address decisions without a central
full-state controller. Our mechanism combines four requirements. It treats
pod subscription count and backlog as separate pressure signals. It tracks
the client's subscription-to-connection assignments. It promotes the same
authenticated probe socket only after checking that candidate. Finally, it
uses receipt-gated delivery handoff while durable queues stay in place.
These requirements
define the comparison in Table~\ref{tab:comparison_related_work}.

The proposed mechanism connects these concerns through an endpoint-preserving
probe workflow. The load balancer selects a reachable Pod when the probe is
opened; the selected policy then evaluates whether to use that already-established connection.
Promotion is conditional, not an automatic consequence of probing. An
admissible migration reuses the observed socket and transfers subscription
delivery ownership under the handoff contract, while the durable queues remain
stationary. This differs from choosing a Pod by name and opening another
load-balanced socket, whose endpoint could differ from the one evaluated.
Local subscription and backlog measurements also support scaling and
redistribution across existing connections. The contribution is this
partially observable control and execution mechanism, evaluated with established
greedy and Q-learning methods. Their comparison is specific to the measured workloads.
Table~\ref{tab:comparison_related_work} summarizes the
positioning of the proposed approach relative to representative literature.

\begin{table*}[t]
	\small
	\caption{Comparison of Representative Work and the Proposed Framework}
	\label{tab:comparison_related_work}
	\centering
	\setlength{\tabcolsep}{3pt}
	\renewcommand{\arraystretch}{1.08}
	\begin{tabular}{>{\raggedright\arraybackslash}p{0.17\textwidth}>{\raggedright\arraybackslash}p{0.18\textwidth}>{\raggedright\arraybackslash}p{0.17\textwidth}>{\raggedright\arraybackslash}p{0.18\textwidth}>{\raggedright\arraybackslash}p{0.18\textwidth}}
		\toprule
		\textbf{Work} & \textbf{Main focus} & \textbf{Observed information} & \textbf{Controlled object} & \textbf{How this work differs}\\
		\midrule
		BLOC \cite{Bhattacharya2024} & Request balancing and overload control & Overload feedback; no central metadata store & Request admission and server selection & Adds retained-socket promotion and subscription ownership transfer\\
		QEdgeProxy \cite{Cilic2025} & QoS-aware service selection & Client QoS outcomes; decentralized bandits & Service-instance selection & Preserves the probed endpoint during connection replacement\\
		Adaptive load balancing \cite{Selvakumar2023} & Service-chain latency & Predicted waiting times; inter-service costs & Request priorities and instance selection & Reassociates existing subscription-bearing connections\\
		Decentralized offloading \cite{Yahya2024} & Task distribution among edge servers & Unknown network information; adjusted server rewards & Offloading population control & Reassociates persistent subscription groups\\
		Per-device DRL \cite{Nieto2024DRLOffloading} & Energy- and latency-aware offloading & Device, server, channel, and task state & Local, MEC, or cloud task execution & Uses subscription counts and queued-event pressure\\
		Partial-observation DRL \cite{Hou2023PODRL} & QoE-aware distributed offloading & Partial observations; device QoE preferences & IoT task-offloading decisions & Evaluates a retained probe and transfers subscriptions\\
		Dr-Stream \cite{Liu2020DrStream} & Stream redirection and operator elasticity & Stream-runtime measurements & Operator instances and stream state & Keeps durable queues stationary during transport handoff\\
		Aj-Stream \cite{Sun2024AjStream} & Skew-aware stream joining & Predicted key frequencies; instance load & Key grouping and operator scaling & Moves subscription ownership while queues stay in place\\
		RL autoscaling \cite{AbdelKhaleq2021Autoscaling} & Microservice scaling thresholds & Resource demand and response-time QoS & Available service capacity & Moves existing sticky sockets to available capacity\\
		Hephaestus\-Forge \cite{Santos2025HephaestusForge} & Multi-cluster microservice placement & Network latency; resource capacity; deployment cost & Pod-replica placement across clusters & Connects existing subscriptions to available Pod capacity\\
		Feature-model-guided RL \cite{Metzger2024GuidedExploration} & Exploration of adaptation choices & Feature models; configuration changes; rewards & Runtime feature configurations & Uses bounded connection actions with receipt-gated handoff\\
		RS-DRL \cite{Kavianifar2026RSDRL} & Reward-shaped self-adaptation & Runtime feedback; replayed transitions & Self-adaptive IoT configurations & Specializes actions to guarded transport and subscription handoff\\
		\bottomrule
	\end{tabular}
\end{table*}


The broader cloud-scheduling and edge-computing context is retained in
Section~\suppref{app:extended-literature}.
